% Chapter 1
% 
\chapter{Introdução} % Main chapter title
\label{chap:Chapter1} % For referencing the chapter elsewhere

%-------------------------------------------------------------------------------
%---------
%
\section{Contextualização} 
\label{sec:chap1_context} %For referencing this section elsewhere

\section{Enquadramento do projeto} 
\label{sec:chap1_enquadramento} %For referencing this section elsewhere

\section{Problema e motivação} 
\label{sec:chap1_problem} %For referencing this section elsewhere

A disponibilização de CBT-I através de aplicações móveis permite apoiar a execução de atividades terapêuticas e a monitorização do sono. Existem soluções digitais personalizadas e estudos que integram wearables; permanece, contudo, relevante investigar como combinar dados individuais de sono e contexto com uma política de notificação ao utilizador adaptativa, e qual o benefício adicional face a uma política predefinida.
O problema deste projeto consiste em conceber e avaliar um mecanismo que selecione recomendações de um percurso terapêutico definido e identifique oportunidades adequadas para as disponibilizar. As decisões devem considerar o histórico e o contexto do utilizador, a qualidade e o instante de disponibilidade dos dados e a possibilidade de adiar ou não enviar uma intervenção. A deteção de necessidade, a disponibilidade do utilizador e o benefício terapêutico constituem dimensões distintas.
Os dados do smartwatch serão complementados por um diário de sono e, quando pertinente, por avaliações breves de contexto. As estimativas dos wearables não serão assumidas como medições clínicas de referência.

\section{Objetivos} 
\label{sec:chap1_objectives} %For referencing this section elsewhere

Objetivo geral: desenvolver e avaliar um protótipo móvel de apoio à CBT-I que utilize dados de smartwatch e informação do utilizador para adaptar recomendações e o seu momento de entrega, investigando a fiabilidade técnica e o possível efeito na adesão e no sono.

Objetivos Específicos:
\begin{itemize}
    \item Identificar os dados de sono e contexto relevantes e verificar a sua disponibilidade, qualidade e atraso de acesso no smartwatch e na plataforma selecionados;
    \item Definir, com apoio de especialistas, o percurso de CBT-I, os conteúdos elegíveis e os critérios para reconhecer oportunidades de intervenção;
    \item Desenvolver e comparar um mecanismo de IA para identificar essas oportunidades com uma referência simples baseada em regras explícitas;
    \item Implementar um sistema de recomendações que adapte a seleção e a entrega de conteúdos ao perfil, histórico e contexto do utilizador, respeitando as regras terapêuticas definidas;
    \item Construir um protótipo de aplicação móvel com diário de sono, recolha de dados, acompanhamento do progresso e notificações adaptativas, integrando o mecanismo de IA;
    \item Avaliar o desempenho técnico, a usabilidade e a viabilidade de utilização; se existirem condições, realizar um piloto comparativo para estudar adesão e alterações no sono;
    \item Analisar resultados, limitações e fatores que condicionam a personalização e o momento de entrega, propondo melhorias para trabalhos futuros.
\end{itemize}


\section{Questões de investigação} 
\label{sec:chap1_researchQuestions} %For referencing this section elsewhere

Pergunta 1:
Qual o efeito de recomendações de CBT-I personalizadas e entregues em momentos adaptados ao contexto do utilizador na qualidade do sono e na adesão ao tratamento, comparativamente a recomendações genéricas entregues em horários fixos?

Pergunta 2:
Com que desempenho é possível identificar oportunidades de intervenção de CBT-I, previamente definidas, utilizando técnicas de inteligência artificial e dados de smartwatch disponíveis no momento da decisão?

\section{Abordagem geral e contributos esperados} 
\label{sec:chap1_approach} %For referencing this section elsewhere

\section{Organização do documento} 
\label{sec:chap1_organization} %For referencing this section elsewhere
